\section{Conclusion}
\label{sec:conclusion}

This paper discussed an automaton-theoretic representation for the
differential distribution tables. We outlined the memory constraints
inherent in differential cryptanalysis, presented a method that reframes
DDTs as formal languages, and proposed representing them using finite
automata to capture structural regularities. We validated this
representation by evaluating its memory usage across various S-boxes,
confirming that the structural compression translates into concrete storage
advantages over standard matrix implementations.

Representing a DDT as a dense matrix provides straightforward access but
suffers from exponential growth, becoming prohibitive for storage as S-box
sizes increase. These constraints force researchers to limit the scope of
their analysis, either by restricting the size of the primitives analyzed
or by relying on simplifications and heuristics. By exploiting the
redundancies present in typical S-boxes, we demonstrated that this tabular
structure can be substantially compressed without any loss of data, opening
the possibility of extending the practical boundaries of analysis.

Beyond the storage gains, the central observation of this work is that the
language-theoretic view does more than reduce memory. Once the DDT is
expressed as a regular language, the cryptanalytic searches that an
analyst performs over the table become operations on automata, and the
wide range of properties that can be phrased as regular conditions over
the quadrant encoding of $(a,b)$ are all recovered through a single
joint traversal of the DDT automaton and a small query automaton. We
described several such queries, covering
fixed-input and truncated differentials, impossible differentials,
extremal-probability entries, low-weight transitions, linearly constrained
pairs, composite queries built by boolean closure, and
cross-DDT comparisons. Each of these tasks, historically implemented by a
bespoke scan of the dense table, becomes a short regular expression in the
proposed framework and inherits the compactness of the underlying
automaton.

Several directions naturally extend the present work. The query catalogue
can be enriched with non-regular but context-free constraints, such as
bounded-difference balances across multiple rounds, by replacing the
query automaton with a pushdown automaton and adapting the joint
traversal accordingly.
The same automaton can also serve as a generator of linear inequalities
for MILP-based trail search, providing an alternative route to the
modeling of S-box behavior; a careful study of the trade-off between
automaton size and the number of generated inequalities is left to future
work. Finally, the size of the reduced automaton itself appears to capture
an intrinsic notion of structural regularity of an S-box, and quantifying
the deviation between an actual DDT and its random or theoretically
optimal counterparts via this size-based metric is a natural analytical
follow-up.
